\documentclass[11pt]{article}
\usepackage[T1]{fontenc}
\usepackage{lmodern}
\usepackage[margin=1in]{geometry}
\usepackage{amsmath,amssymb,amsthm,booktabs}
\usepackage{xurl}
\usepackage[hidelinks]{hyperref}
\clubpenalty=10000
\widowpenalty=10000
\newtheorem{theorem}{Theorem}[section]
\newtheorem{proposition}[theorem]{Proposition}
\newtheorem{lemma}[theorem]{Lemma}
\newtheorem{corollary}[theorem]{Corollary}
\theoremstyle{definition}
\newtheorem{definition}[theorem]{Definition}
\newtheorem{remark}[theorem]{Remark}
\newcommand{\NN}{\mathbb N}
\newcommand{\ZZ}{\mathbb Z}
\newcommand{\QQ}{\mathbb Q}
\newcommand{\A}{\mathcal A}
\title{Atoms and Non-atomicity in the Rational Semiring\\
$\mathbb N_0[2/5,2/3]$}
\author{Rational Sumring Lab\\\small Working draft; author attribution to be reviewed}
\date{November 1, 2026}
\begin{document}
\maketitle
\begin{abstract}
Let $S=\mathbb N_0[2/5,2/3]$. We prove that its nonzero multiplicative
monoid is not atomic: the element $3$ has no finite factorization into
atoms. The proof uses a binary digit map, the sharp real threshold $5/3$
above which every element of $\mathbb Z[1/15]$ belongs to $S$, and a
restriction on the denominators of atoms. We also describe all atoms in
terms of binary polynomials and eventually constant binary expansions,
and give a terminating exact membership procedure. The IDF property is
reduced to an explicit finiteness question for irreducible binary
polynomials and is not decided here. Non-atomicity of this exact semiring
was announced in a June 2026 CrowdMath post without a written proof;
we provide a complete argument and credit earlier digit and perturbation
methods. Priority of the additional results has not been established.
\end{abstract}

\section{Statement and provenance}
Throughout, set
\[
 q=\frac25,\qquad r=\frac23,\qquad h=\frac53=1+r=\frac1{1-q},
 \qquad R=\ZZ[1/15],\qquad S=\NN_0[q,r],\qquad M=S\setminus\{0\}.
\]
The operation on $M$ is multiplication. A unit has its reciprocal in $M$.
An atom is a nonunit which cannot be expressed as a product of two
nonunits. A monoid is atomic if every nonunit is a finite product of
atoms. It is irreducible-divisor finite (IDF) if every element has only
finitely many atom divisors up to multiplication by units.

\begin{theorem}\label{thm:main}
The element $3\in M$ has no finite factorization into atoms of $M$.
Consequently, $M$ is not atomic.
\end{theorem}

This addresses the atomicity part of CrowdMath 2026 Open Problem 2(d)
\cite{cm2}. Its IDF part remains unresolved in this note. In particular,
non-atomicity is not used as a reason to infer IDF or non-IDF.

\paragraph{Prior work and attribution.}
The exact non-atomicity conclusion was announced by the participant
BakedPotatoz in progress-log post~7 on June 27, 2026, with the proof
expressly left unwritten~\cite{announcement}. Thus we do not claim first
discovery of that conclusion. Digit expansions and perturbations by
technical fractions near one occur in the earlier project writeups
\cite[Section 4.6]{progress} and \cite[Theorem 2.6]{yayyayyay}.
The latter also proves valuation bounds for two generators
\cite[Proposition 3.7]{yayyayyay}. We adapt these methods and give every
argument needed here, including the restriction on the power of $3$ in
an atom's denominator and the obstruction to a factorization of $3$.
A complete earlier proof of this exact conclusion was not located in
our limited source audit; that does not establish priority. The
one-generator results of Deng, Gotti, and Zeng~\cite{dgz} supply related
context but are not invoked to transfer atomicity to this two-generator
semiring.

This draft was prepared with AI assistance. Its complete written proofs
and executable consistency checks have undergone independent agent
review, not human peer review or proof-assistant verification. Author
attribution, priority, and shared CrowdMath contributions require review
before submission. No MIT affiliation or endorsement is claimed.

\section{Parity, a prime element, and binary digits}
Every element of $S$ lies in $R$, and every nonzero element has
nonnegative valuation at every prime outside $\{3,5\}$. In particular,
$v_2(x)\ge0$ for $x\in M$. The identities
\begin{equation}\label{eq:identities}
 r=q+qr=qh,\qquad h=1+r
\end{equation}
will be used repeatedly.
For any $x\in R$, let $\epsilon(x)\in\{0,1\}$ be its residue modulo
$2R$, and define
\[
 \tau(x)=\frac{x-\epsilon(x)}q.
\]
This lies in $R$: denominators in $R$ are odd, and division of
$x-\epsilon(x)\in2R$ by $q=2/5$ introduces no new denominator prime.

\begin{lemma}\label{lem:tau}
The map $\tau$ sends $S$ into $S$. If $x\in R$ and $\tau(x)\in S$,
then $x\in S$. For $x\in M$, divisibility by $q$ inside $M$ is equivalent
to $v_2(x)>0$.
\end{lemma}
\begin{proof}
Choose a polynomial representation of $x\in S$ with constant coefficient
$c$. Every nonconstant monomial has positive $2$-valuation, so
$c\equiv\epsilon(x)\pmod2$. The nonnegative even integer
$c-\epsilon(x)$ divided by $q$ is a nonnegative integer. Every
nonconstant monomial is divisible by $q$ inside $S$: this is immediate
if its $q$-exponent is positive, and otherwise follows from $r/q=h\in S$.
Dividing the representation of $x-\epsilon(x)$ by $q$ proves
$\tau(x)\in S$. The converse follows from
$x=\epsilon(x)+q\tau(x)$. If $v_2(x)>0$, then $\epsilon(x)=0$ and
$x/q=\tau(x)\in M$. Conversely, $x=qy$ with $y\in M$ gives
$v_2(x)=1+v_2(y)>0$.
\end{proof}

\begin{proposition}\label{prop:split}
The only unit of $M$ is $1$. The element $q$ is prime, and, writing
$T=\{x\in M:v_2(x)=0\}$, the map
\[
 \NN_0\times T\longrightarrow M,\qquad (k,t)\longmapsto q^kt
\]
is a monoid isomorphism, with addition in the first coordinate.
Every element of $T$ is at least one. In particular, the only atom
of positive $2$-valuation is $q$.
\end{proposition}
\begin{proof}
If $x\in T$, the constant coefficient in any representation is odd,
so it is at least one and $x\ge1$. A nontrivial unit would have either
itself or its inverse below one. That element has positive $2$-valuation,
so its inverse has negative valuation and cannot lie in $M$.
The divisibility criterion in Lemma~\ref{lem:tau} and additivity of
$v_2$ show that $q\mid ab$ implies $q\mid a$ or $q\mid b$.
Repeated division by $q$ gives $x=q^kt$ with $k=v_2(x)$ and $t\in T$;
these values determine both factors uniquely. This proves the
isomorphism. Since $q$ is a nonunit prime, it is an atom, and every
other element of positive valuation has a nontrivial factorization
with a factor $q$.
\end{proof}

Set $c_i(x)=\epsilon(\tau^i(x))$ and
$\mu_k(x)=\sum_{i=0}^{k-1}c_i(x)q^i$. Iteration gives
\begin{equation}\label{eq:prefix}
 x=\mu_k(x)+q^k\tau^k(x).
\end{equation}
If $x-y\in2^kR$, their first $k$ digits agree and
\begin{equation}\label{eq:congruence}
 \tau^k(x)-\tau^k(y)=q^{-k}(x-y).
\end{equation}
Indeed, their parities initially agree, and division by $q$ lowers
the guaranteed $2$-divisibility by one; induction proves both statements.
The identities hold for every $x,y\in R$, including negative iterates.

\section{A sharp membership threshold}
\begin{lemma}\label{lem:coarse}
Every $x\in R$ with $x\ge20/3$ belongs to $S$.
\end{lemma}
\begin{proof}
Write the denominator of $x$ as $3^a5^b$. The Chinese remainder theorem
splits its class modulo $\ZZ$ into a class with denominator $5^b$ and
one with denominator $3^a$.
Every class with denominator $5^b$ has a representative
$u=\sum_{i=1}^b e_iq^i$ with $0\le e_i\le4$: at the highest degree,
choose $e_b$ so that $e_b2^b$ matches the target numerator modulo $5$,
subtract that term, and divide the remaining numerator by $5$.
The same induction at the lower degrees finishes the construction.
Similarly, every class with denominator $3^a$ has a representative
$v=\sum_{j=1}^a f_jr^j$ with $0\le f_j\le2$.
These nonnegative representatives lie in $S$ and satisfy
\[
 u<4\sum_{i\ge1}q^i=\frac83,
 \qquad v<2\sum_{j\ge1}r^j=4.
\]
Consequently $x-u-v$ is a nonnegative integer when $x\ge20/3$,
which proves $x\in S$.
\end{proof}

\begin{theorem}\label{thm:conductor}
Every $x\in R$ with $x\ge h=5/3$ lies in $S$.
No smaller real threshold has this property.
\end{theorem}
\begin{proof}
Since $h(1-q)=1$ and $\epsilon(x)\le1$,
\begin{equation}\label{eq:growth}
 \tau(x)-h=q^{-1}\bigl(x-h+1-\epsilon(x)\bigr)
 \ge q^{-1}(x-h).
\end{equation}
For $x>h$, its iterates stay above $h$ and eventually exceed $20/3$.
Lemma~\ref{lem:coarse} places that iterate in $S$, and repeated use of
Lemma~\ref{lem:tau} places $x$ in $S$. The endpoint $h=1+r$ also lies
in $S$.

For sharpness let $x_n=h-2q^n$ for $n\ge1$. These are positive
members of $R$, increase to $h$, and differ from $h$ by an element
of $2^{n+1}R$. Since $\tau(h)=h$, equation~\eqref{eq:congruence} gives
$\tau^n(x_n)=h-2=-1/3$. Such an iterate is impossible for an element
of $S$, by Lemma~\ref{lem:tau}. Thus $x_n\notin S$.
\end{proof}

Let $G=\{3^a5^b:a,b\in\ZZ\}$, the group of positive units of $R$.
Its intersection with $M$ is the monoid of technical fractions; it is
smaller than $T$, which only requires zero $2$-valuation.

\begin{lemma}\label{lem:small}
For every $L\in\NN_0$ and $\eta>0$, there is $y\in M\cap G$ with
$1<y<1+\eta$ and $y\equiv1\pmod{2^LR}$.
\end{lemma}
\begin{proof}
Choose $j\ge\max(1,L)$ with $hq^j<\eta$.
The subgroup generated by $3^{2^j}$ and $5^{2^j}$ is dense in
$\mathbb R_{>0}$. To see this, unique prime factorization makes
$\log3/\log5$ irrational; integer combinations of these logarithms
are dense in $\mathbb R$, and multiplication by $2^j$ preserves
this density. Equivalently, arbitrarily small nonzero combinations
follow from the pigeonhole principle, and their integer multiples
approximate any real number.
Choose $y$ in that subgroup satisfying
\[
 1+hq^j<y<1+\eta.
\]
Euler's theorem for odd residues modulo $2^j$ shows that the subgroup
is contained in $1+2^jR$, including its negative powers.
Thus $z=(y-1)/q^j$ lies in $R$ and exceeds $h$.
Theorem~\ref{thm:conductor} gives $z\in S$, and
$y=1+q^jz\in M\cap G$ has all the required properties.
\end{proof}

\section{Restricting the atoms}
Define the set
\[
 \mathcal B=\{x\in S:0\le\tau^k(x)\le h\text{ for every }k\in\NN_0\}.
\]
It is essential here that membership in $S$ is part of this definition.

\begin{lemma}\label{lem:atomsboundary}
Every atom other than $q$ belongs to $T\cap\mathcal B$ and exceeds one.
\end{lemma}
\begin{proof}
Proposition~\ref{prop:split} gives membership in $T$ and $x>1$.
Suppose $\tau^k(x)>h$ for some $k\ge0$.
Choose $\eta>0$ so small that $1+\eta<x$ and
\[
 x\left(1-\frac1{1+\eta}\right)
 <q^k\bigl(\tau^k(x)-h\bigr).
\]
Lemma~\ref{lem:small} supplies a technical fraction $y$ with
$1<y<1+\eta$ and $y\equiv1\pmod{2^kR}$.
Because $y$ is a unit of $R$, $x/y\in R$ and $x/y\equiv x\pmod{2^kR}$.
Equation~\eqref{eq:congruence} therefore gives
\[
 \tau^k(x/y)=\tau^k(x)-\frac{x-x/y}{q^k}>h.
\]
The conductor theorem and the backward implication of
Lemma~\ref{lem:tau} give $x/y\in S$.
Both $y$ and $x/y$ exceed one, contradicting the assumption that $x$
is an atom.
\end{proof}

\begin{lemma}\label{lem:nooverlap}
No element $x\in\mathcal B$ can be written as $2q^n+z$ with
$n\in\NN_0$ and $z\in S$.
\end{lemma}
\begin{proof}
Since $x-z=2q^n\in2^{n+1}R$, their first $n$ digits agree. Hence
$\tau^n(x)=2+\tau^n(z)\ge2>h$, a contradiction.
\end{proof}

\begin{proposition}\label{prop:forms}
An element $x\in S$ admits no expression $2q^n+z$ as in
Lemma~\ref{lem:nooverlap} if and only if it has one of the forms
\begin{align}
 x&=P(q),\label{eq:finite}\\
 x&=P(q)+q^nr,\quad n\ge0,\quad\deg P\le n,\label{eq:tail}
\end{align}
where $P$ has coefficients in $\{0,1\}$ (and the zero polynomial is
allowed). These elements are exactly $\mathcal B$.
An element of the second form has $v_3(x)=-1$.
\end{proposition}
\begin{proof}
Choose any nonnegative polynomial representation in $q,r$.
A term $q^ir^j$ with $j\ge2$ yields a summand $2q^{i+j+1}$:
expand $r^j=q^j(1+r)^j$, retain the term $j q^jr$, and use
$r=q+qr$. All discarded terms remain in $S$.
Consequently every $r$-exponent is at most one.
Two linear $r$-terms, including two copies of the same term,
also yield an overlap, because for $N>i$,
\begin{equation}\label{eq:tailidentity}
 q^ir=\sum_{m=i+1}^{N}q^m+q^Nr.
\end{equation}
Choosing $N$ large enough exposes a common power of $q$ in both terms.
Thus there is at most one linear $r$-term, with coefficient one.
Pure $q$-terms have coefficients at most one; if the linear term is
$q^nr$, no pure term can have degree greater than $n$, again by
\eqref{eq:tailidentity}. This proves the asserted forms.

Conversely, their digit iterates remain in $[0,h]$. A finite binary
polynomial loses its constant digit at each application of $\tau$
and eventually reaches zero. For \eqref{eq:tail}, after $n$ steps
one reaches either $r$ or $1+r=h$; moreover $\tau(r)=h$ and
$\tau(h)=h$. Each preceding binary prefix with its remaining tail
is at most $\sum_{i\ge0}q^i=h$. Thus both forms belong to
$\mathcal B$, and Lemma~\ref{lem:nooverlap} excludes an expression
$2q^n+z$.
If $P=0$, the final assertion is immediate. Otherwise,
$v_3(P(q))\ge0$ while $v_3(q^nr)=-1$, so the latter valuation
cannot cancel and $v_3(x)=-1$.
\end{proof}

\begin{corollary}\label{cor:denominator}
Every atom other than $q$ lies in $(1,5/3]$ and has reduced denominator
with exponent at most one at the prime $3$.
\end{corollary}

\section{Why the element 3 has no atom factorization}
\begin{proof}[Proof of Theorem~\ref{thm:main}]
If $3$ were a product of atoms, every factor would be technical.
Indeed, valuations outside $\{3,5\}$ are nonnegative on $M$ and sum
to zero in a factorization of $3$. No numerator prime outside
$\{3,5\}$ can therefore occur in any factor. In particular $q$
cannot be a factor.

Let $a$ be any technical atom. By Corollary~\ref{cor:denominator},
$1<a\le5/3$ and $v_3(a)\ge-1$.
If $v_3(a)=-1$, technicality gives $a=5^m/3$ for some integer $m$.
The interval restriction forces $m=1$, so $a=h$.
Otherwise $a=3^A/5^B$ with integers $A>B\ge1$: nonnegative
$5$-valuation would make $a$ an integer greater than one, hence at
least $3$; and $a>1$ forces $A>B$.

Consider the additive valuation weight $\delta=v_3+v_5$ on $G$.
The candidate $h$ has $\delta(h)=0$; every other technical atom has
$\delta(a)=A-B\ge1$. Equality to one would require
\[
 a=\frac{3^{B+1}}{5^B}=3(3/5)^B.
\]
For $B=1$ this is $9/5>5/3$; for $B\ge3$ it is at most
$81/125<1$. The only remaining value is $27/25$ at $B=2$.
But
\[
 \tau(27/25)=1/5,\qquad \tau^2(27/25)=-2,
\]
so $27/25\notin S$ by Lemma~\ref{lem:tau}.
Thus a technical atom has weight either zero or at least two.
No finite sum of these weights can equal $\delta(3)=1$.
This rules out every finite atom factorization of $3$.
\end{proof}

The proof does not rely on an exhaustive search for atoms. The restriction
on their denominators applies to every representation of every possible
atom, and the final weight argument excludes arbitrarily long products.

\section{A complete description of the atoms}
A binary polynomial is a polynomial whose coefficients are zero or one.
Irreducibility below is in the polynomial semiring $\NN_0[X]$, whose
only multiplicative unit is the constant polynomial $1$.

\begin{theorem}\label{thm:classification}
The atoms of $M$ are exactly:
\begin{enumerate}
\item $q$;
\item $P(q)$, where $P$ is a nonconstant binary polynomial with $P(0)=1$
which is irreducible in $\NN_0[X]$;
\item $P(q)+q^nr$, where $n\ge0$, $P$ is binary, $P(0)=1$, and
$\deg P\le n$.
\end{enumerate}
Repetitions within the third description are harmless. Its canonical
form is $h$, or $n\ge1$ with $\deg P<n$.
\end{theorem}
\begin{proof}
Lemmas~\ref{lem:atomsboundary} and \ref{lem:nooverlap}, together with
Proposition~\ref{prop:forms}, prove that every atom other than $q$
has one of the two displayed forms with constant digit one.

Suppose $x\in\mathcal B\cap T$ has a factorization $x=ab$ with both
factors different from one. Additivity of $v_2$ gives $v_2(a)=v_2(b)=0$.
If a factor contained a summand $2q^k$,
choose any monomial in a representation of the other factor and
expand $r=q(1+r)$ to extract a pure power of $q$. Their product would
contain $2q^N$, contrary to Lemma~\ref{lem:nooverlap}.
Thus both factors have the forms of Proposition~\ref{prop:forms},
with constant digit one.
If one factor has a linear $r$-term, its product with the other's
constant term supplies one tail, and its product with any nonconstant
term of the other supplies another tail or an $r^2$-term. The finite
expansions in the proof of Proposition~\ref{prop:forms} then produce
$2q^N$, again a contradiction. Therefore both factors are finite
binary polynomials evaluated at $q$.

This already proves that every element in item 3 is an atom:
it has $3$ in its denominator, while a product of two finite
polynomial values has denominator supported on $5$ alone.
For item 2, write the factors as $F(q),G(q)$. Their polynomial
product has binary coefficients, since any coefficient at least two
would give a forbidden summand $2q^N$.
Evaluation at $q<1/2$ is injective on binary polynomials: at the
first differing degree $k$, the absolute contribution $q^k$
exceeds the sum of all possible later contributions,
$q^{k+1}/(1-q)<q^k$.
Hence $P(q)=F(q)G(q)$ implies the polynomial identity $P=FG$.
Conversely, a nontrivial factorization of a binary polynomial
$P$ with constant term one in $\NN_0[X]$ necessarily has binary
factors with constant term one, and their values at $q$ exceed one.
It gives a nontrivial factorization in $M$.
Finally, if the coefficient of $X^n$ in the tail form is one,
$q^n+q^nr=q^{n-1}r$ reduces $n$ when $n\ge1$.
Repeatedly doing so gives the canonical forms stated.
\end{proof}

\begin{corollary}
The elements $h=5/3$, $1+q^n$ for $n\ge1$, and $1+q^nr$ for
$n\ge0$ are atoms. There are infinitely many atoms tending to one
from above.
\end{corollary}
\begin{proof}
The tail family follows from item 3. If $1+X^n=FG$ is a nontrivial
factorization in $\NN_0[X]$, then $F(0)=G(0)=1$ and
$0<\deg F<n$. The leading term of $F$ times $G(0)$ forces a positive
coefficient in degree $\deg F$, a contradiction. Item 2 applies.
The limits and distinctness follow from $0<q<1$.
\end{proof}

\section{Exact membership and the remaining IDF question}
\begin{theorem}\label{thm:membership}
Membership in $S$ is decided by the following exact procedure.
Reject a rational outside $R$ or below zero. For $x\in R_{\ge0}$,
iterate $\tau$, accepting upon reaching $0$ or a value at least $h$,
and rejecting upon reaching a negative value or repeating a previous
value without having accepted. This procedure terminates.
\end{theorem}
\begin{proof}
Acceptance is justified by Theorem~\ref{thm:conductor} and repeated
backward use of Lemma~\ref{lem:tau}. A negative iterate is impossible
for a member of $S$.
If a cycle is reached before acceptance and the original $x$ belonged
to $S$, all its iterates would be in $[0,h]$, so
Proposition~\ref{prop:forms} would apply. Those forms reach either
$0$ or $h$, as its proof shows, contradicting the detected cycle.

Write the initial reduced denominator as $3^a5^b$.
Each application of $\tau$ decreases the exponent of $5$ by one
when it is positive, and never increases the exponent of $3$.
After at most $b$ steps the iterates have denominators dividing $3^a$.
Unless the procedure stops earlier, they stay in the bounded interval
$[0,h)$, which contains only finitely many points of $3^{-a}\ZZ$.
A repeat must therefore occur. This proves termination, not a
polynomial-time complexity bound.
\end{proof}

For an atom $a\ne q$, let
\[
 d(a)=\prod_{p\notin\{3,5\}}p^{v_p(a)}.
\]
This is a positive integer coprime to $30$; its $2$-valuation is zero.

\begin{proposition}\label{prop:idf}
For every integer $d\ge1$ coprime to $30$, let
$\A_d=\{a\in\A(M)\setminus\{q\}:d(a)=d\}$.
Every member of $\A_d$ divides $3d$ in $M$.
The monoid $M$ is IDF if and only if each set $\A_d$ is finite.
\end{proposition}
\begin{proof}
For $a\in\A_d$, the rational $3d/a$ belongs to $R$, because all
numerator primes of $a$ outside $\{3,5\}$ are cancelled by $d$.
Also $a\le h$, so $3d/a\ge3/h=9/5>h$.
The conductor theorem puts $3d/a$ in $M$.
If $M$ is IDF, trivial units and this divisibility force each
$\A_d$ to be finite.
Conversely, for fixed $x\in M$, an atom divisor other than $q$ has
$d(a)$ dividing the fixed positive integer
$\prod_{p\notin\{3,5\}}p^{v_p(x)}$. There are only finitely many
such divisors $d$, and a finite union of the corresponding finite
sets, together with $q$, contains all atom divisors of $x$.
\end{proof}

The tail atoms in any fixed $\A_d$ are finite in number.
For a canonical tail with $n\ge1$, its reduced numerator is
\[
 N=3\sum_{k=0}^{n-1}\epsilon_k2^k5^{n-k}+2^{n+1}.
\]
It is odd and divisible by neither $3$ nor $5$, so $d=N$.
Since $1<a\le h$, it satisfies $3\cdot5^n<d\le5\cdot5^n$.
These intervals are disjoint for different $n$, and binary evaluation
is injective for fixed $n$. Thus there is at most one such tail for
a given $d$, in addition to the separately listed atom $h$ with $d=1$.

Consequently, IDF is equivalent to the following precise finiteness
statement: for every fixed $d$ coprime to $30$, there are only finitely
many nonconstant irreducible binary polynomials $P\in\NN_0[X]$ with
$P(0)=1$ such that
\begin{equation}\label{eq:idfremaining}
 5^{\deg P}P(2/5)=d\,3^A\qquad\text{for some }A\in\NN_0.
\end{equation}
Indeed the numerator on the left is odd and not divisible by $5$,
so this is exactly the finite-polynomial part of the fiber condition.
No proof or counterexample to this finiteness statement is supplied
here. Infinite atom families alone do not disprove IDF: their numerator
prime factors can change, so they need not divide one fixed element.

\section{Implementation and checks}
The companion implementation uses integers and exact rational arithmetic.
The module \path{dyadic_semiring.py} implements the parity map and the
terminating membership procedure. Positive decisions supply polynomial
certificates; negative decisions supply a domain obstruction or an exact
rejection trajectory. A separate verifier checks the supplied evidence.
The conductor construction supplies explicit nonnegative sums in $q,r$.
The finite checks in \path{paper/verify_non_atomicity.py} verify
identities, prefix congruences, certificate values, and the exceptional
candidate $27/25$. They also check representative parameters of the
proved families. The general non-atomicity and classification theorems
are established by the written arguments, not by finite testing.

The earlier effective-prime algorithms and their proofs are retained
in \path{paper/effective_primes_note.tex}. Their scope is separate
from the atom classification proved here. Regression tests continue
to check the other two requested rational pairs. The broader mixed-pair
atomicity and IDF problems are not resolved by this paper.

\begin{thebibliography}{9}
\raggedright
\bibitem{cm2}
MIT PRIMES / Art of Problem Solving CrowdMath,
\emph{Arithmetic Conjectures on Rational Semirings, Open Problem 2},
updated October 15, 2026; consulted November 1, 2026.
\url{https://artofproblemsolving.com/polymath/mitprimes2026/p/open-prob2}.
\bibitem{announcement}
BakedPotatoz (CrowdMath participant),
\emph{Crowdmath 2026 Progress Log}, post 7, June 27, 2026.
\url{https://artofproblemsolving.com/polymath/mitprimes2026/f/c4514724h3837652_crowdmath_2026_progress_log}.
\bibitem{progress}
CrowdMath participants, \emph{Crowdmath 2026 Progress Log},
consolidated forum draft, consulted November 1, 2026, Section 4.6.
\url{https://cdn.artofproblemsolving.com/attachments/2/6/b981b69621745784dbd3f49ce0a5111f1b8352.pdf}.
\bibitem{yayyayyay}
yayyayyay (CrowdMath participant), \emph{CrowdMath 2026 Progress},
forum draft, consulted November 1, 2026, Theorem 2.6 and Proposition 3.7.
\url{https://cdn.artofproblemsolving.com/attachments/9/3/c0de521bed56113ef73180b261b869245fcd94.pdf}.
\bibitem{dgz}
Anna Deng, Felix Gotti, and Jason Zeng,
\emph{Factorizations in rational monogenic semidomains},
arXiv:2607.10178 (2026).
\url{https://arxiv.org/abs/2607.10178}.
\end{thebibliography}
\end{document}
